% AI-written, AI cross-reviewed; no human sign-off; not a claim that RH is proved
% RH Notes blueprint source excerpts for Z3.2
% Verbatim excerpts in the notation of the parent blueprint. The labelled editions and source ranges appear below.
% See the unit YAML for the current proof status. Earlier-round passages retain their original wording.

% Start of atlas edition annotation. The following block is outside the original excerpt.
\paragraph{Status of named inputs in this reading.}
\noindent\nolinkurl{Formalism.MainNotions.TruncatedSeriesCountingIdentities.admBK} --- hypothesis at this head (\path{Formalism/SharpRestriction/LevelLetterCocycles/DecompositionData.lean:238--243}); not compared.\par
\noindent\nolinkurl{Formalism.WeaklyAdmissibleFiltrations.BlochKatoDim} --- hypothesis at this head (\path{Formalism/WeaklyAdmissibleFiltrations/Statement.lean:342--349}); not compared.\par
External theorem numbers in the original excerpt are not verified here unless a separate edition-specific source record documents the comparison.
% End of atlas edition annotation.

\section*{Z3.2. Local comparisons and global lattice index}

% Parent source SHA-256: 1af64922b98ed2310744411b1f18b25cc426c97220646321471836e1a3a409d6
% blueprint/ch_P.tex lines 9024 through 9335; sec:P-r4-Z32
\subsection{The local arithmetic comparisons in Z3.2}\label{sec:P-r4-Z32}

\paragraph{The first block at a row.}
Fix $w\in R$. Use the coefficient isomorphism for the self extension
and the homeomorphism from the absolute local group onto $D_w$.
Write $T_1$ for the integral source selected by the rank-one
idempotent, $T'_w$ for the integral target, and $f:T_1\to T'_w$ for the
surjective projection followed by the row-coordinate isomorphism.
Its rational extension intertwines the actions. To see the source
among the representations allowed by the row hypothesis, embed its
localization in the evaluated matrices. Matrix rows separate points.
Right multiplication by the inverse family matrix acts on these rows
by the dual representation, with the prescribed character and cyclotomic
twist. In a finite-dimensional space, finitely many of the row evaluations
already have zero common kernel: as long as the kernel is nonzero, select
an evaluation nonzero on it and decrease its dimension.
The normalized target in block coordinates is the prescribed matrix
lattice. For the two lattice embeddings $j_1,j_2$ and the rational
map $f_K$ one has $f_Kj_1=j_2f$. Given an element of the transported
target lattice, write it as the transport of $j_2(t)$, choose $s$
with $f(s)=t$, and transport $j_1(s)$. This lifts that element of
the target lattice (\nolinkurl{latHom_onto},
\path{Formalism/AdmissibleDesigns/FinDimOfLat.lean:173--181};
\nolinkurl{rowTie_disp},
\path{Formalism/AdmissibleDesigns/FiniteClassesRowEqGeneric.lean:109--142}).
Continuous family orbits give continuous orbits on the source. Lift a
target vector and apply the continuous rational projection to its orbit
to obtain target continuity
(\path{rowSub_at}, \path{IdempotentPartRepresentationAt.lean:135--217};
\path{rowMat_SN}, \path{PhiWRhoT.lean:160};
\path{rowSub_S}, \path{rowOrb_S}, \path{rowFin_S},
\path{RowLatticeFiniteFree.lean:62--122};
\path{rowTie_disp}, \path{FiniteClassesRowEqGeneric.lean:102--158}).

An integral finite cocycle is one whose image in the rational
representation becomes a period coboundary. The induced map on period
coefficients takes a witness for the source to one for the target.
It also takes ordinary coboundaries to ordinary coboundaries. Thus $f$
maps finite classes into finite classes. Conversely the definitions of
both finite conditions can be read in the local period model by pulling
a cocycle back along the inverse group and coefficient isomorphisms.
Continuity is preserved in both directions by the lattice and vector
space topologies. More explicitly, let $\iota$ identify the local group with $D_w$ and
let $e$ denote the coefficient identification. On target cochains send
$c$ to $g\mapsto e(j_2c(\iota g))$. It is injective since $\iota$ is
onto and $j_2$ is injective. A local finite cocycle valued in the
transported lattice pulls back pointwise through $j_2$ and
$\iota^{-1}$. The topology of the embedded lattice makes this
pullback continuous. Equivariance gives the cocycle identity, and the
period test transfers by the same coefficient transport. Rational
coboundaries belong to the finite condition, so the criterion descends
to integral classes. For the denominator, the image cocycles on the
local side are sums of the projected finite source cocycles and
integral target coboundaries. Transport sends both summands to the
corresponding summands. Conversely lift a source representative by
the source lattice inverse and pull back the integral primitive of
the second summand. This identifies the inverse image of the
denominator with $J_w$ on classes, including the coboundary
ambiguity. The resulting semilinear bijection on quotients identifies
their submodule chains, hence their lengths
(\nolinkurl{exists_pre}, \nolinkurl{im_iff}, \nolinkurl{len_eq_col},
\path{Formalism/AdmissibleDesigns/LiftMapEquivariant.lean:164--305}). Apply the assumed
\path{RowAt} bound to these two lattices and this particular map:
it gives integers $m,k$ with colength $m$, $k=\operatorname{length}_O(O/lO)$,
and $m\le b_w k$. Here $k=a$. This proves both $J_w\subset F_{s,w}$
and finiteness with the required bound
(\path{RowAt}, \path{Formalism/SharpRestriction/LevelLetterCocycles/QuotientHomNonvanishingBound.lean:142--147};
\path{im_le_Fs}, \path{FiniteClassesContinuousLattice.lean:560};
\path{fs_eq}, \path{finite_classes_row_eq_generic}, \path{rowLenD_eq},
\nolinkurl{rowTie_disp},
\path{Formalism/AdmissibleDesigns/FiniteClassesRowEqGeneric.lean:40--100,102--157}).

\paragraph{The cup-product test.}
For degree-one cocycles $u,v$, the cup representative at $(g,h)$
is $b(u(g),g v(h))$, with the given pairing $b$ into the cyclotomic
twist. Under the same local identifications this representative and the
coboundary of any continuous one-cochain are carried to the corresponding
representatives in the period model. Indeed the coefficient map commutes
with addition, multiplication and the cyclotomic action. Transporting
back proves the converse coboundary implication. The public annihilator
hypothesis for the de Rham representation and its paired dual now
identifies the rational finite condition with the annihilator of the
dual rational finite condition. The local invariant isomorphisms on
degree two preserve zero. Write $F,F'$ for the two rational finite spaces, $H'_{T}$ for
the dual integral cohomology and $j':H'_T\to H'_{V}$ for its
rationalization. Write $q$ for the reduction of rational classes on
the first side and $U=K/O$. The duality compatibility gives
$\langle qx,t\rangle_U=\langle x,j't\rangle_K\bmod O$.
Thus $q(F)$ annihilates $F_{s,w}=(j')^{-1}F'$.

For the converse take a torsion-side class $z$ annihilating
$F_{s,w}$. Its functional $t\mapsto\langle z,t\rangle_U$
factors through the image of $H'_T$ in $H'_V/F'$. This image is
finitely generated by integral Tate finiteness and torsion-free
inside a $K$-space, hence free over the DVR $O$. Choose an $O$-basis,
lift its functional values from $U$ to $K$, and extend the resulting
functional $K$-linearly to $H'_V/F'$. The basis stays independent
after inverting the uniformizer. Rational duality represents the
composite on $H'_V$ by a class $x$. Since that functional vanishes
on $F'$, rational orthogonality gives $x\in F$. The reduction of
this functional agrees on every basis element and therefore on
every integral class with the functional of $z$. Integral
nondegeneracy now gives $qx=z$. This proves both inclusions
(\nolinkurl{bkbConv},
\path{Formalism/WachLatticeTwists/CharacterSplitJumps/CharacterJumpProofs/PropLocBK.lean:13--84}). These are applications of the local
duality hypotheses in the displayed theorem
(\path{cup_iff}, \path{bk_a_of_admitted_bk_a},
\path{CoboundaryMemFiniteClasses.lean:397--440,518--530};
\path{bkInIso}, \path{bkInRow}, \path{BkInRow.lean:14--41};
\path{rowRhoReal}, \path{RowChiReal.lean:16--53}).

There is also a cochain calculation particular to $J_w$.
In the left $A$-submodule $Y=\sum_g A y_g$ of $A$, with $y_g=(u(g)-1)P$, take the coefficient
pairing $T_1\times Y\to Y(1)$ given by multiplication with the
image of the first argument in $A$. Equivariance follows by cancelling
$r(g)^{-1}r(g)$ and the two character factors. In $A$ the cocycle
$g\mapsto y_g$ is the coboundary of $P$. Hence the cup with an
integral source cocycle is a coboundary after inclusion in $A(1)$
(the cup-coboundary identity follows by expanding the two cocycle
identities). It is already a continuous coboundary in $Y(1)$:
the local invariant isomorphisms for finitely generated $O$-modules with
their module topologies identify the map on $H^2$ for the twist with
the inclusion $Y\hookrightarrow A$. Naturality and injectivity of
that inclusion prove injectivity on $H^2$; $Y$ is finite over $O$ as
a submodule of the finite $O$-module $A$. Applying the functional
induced by $\phi$ sends the resulting primitive to a primitive for the
row pairing with $\Psi_0(\phi)$. Thus the pairing vanishes on $J_w$.
Together with the preceding paragraph this gives precisely the
injection of the row error into $(F_{s,w}/J_w)^\vee$ used earlier
(\nolinkurl{Ymod}, \path{Formalism/CharacterColumnKernels/StatementN.lean:92--106}; \path{clause4GK}, \path{Clause4GK.lean:17--20};
\path{h2Map_injective_statement}, \path{clause4},
\path{H2MapInjectiveStatement.lean:14--55};
\path{cupKillAdm}, \path{CupKillAdm.lean:254--274};
\path{rowCup}, \path{KeyMul.lean:12--59};
\path{bridgeRho}, \path{Formalism/WachLatticeTwists/CoprimeLiftBounds/CoprimeLiftProofs/Bridge.lean:14--43}).

\paragraph{Character blocks at the same row.}
Write $e=c-1$, $1\le c\le t$. The coefficient isomorphism identifies
the torsion block, with its action, with $C_e=V_e/L_e$ for the
one-dimensional twist $\epsilon^{-e}$. On $Y\oplus O$ use the extension
action with invariant second coordinate $p$ and base vector $s=(0,1)$.
Then $gs-s=(y_g,0)$. Its rational representation embeds into a finite
sum of the twisted family representations. Concretely, in the
evaluated algebra the map $(y,a)\mapsto y+aP$ intertwines the
action. Record $a$ also on a fixed nonzero vector in the twisted
special fibre (the first character line). These two records are jointly
injective. Finite evaluation coordinates suffice by the decreasing
kernel argument above. The range of the resulting embedding is stable, so the extension
and the quotient fit in an equivariant short exact sequence with
middle term this finite direct sum. Induction on the number of
summands, using the split exact sequence for the last summand and
the Fontaine hypotheses, gives a crystalline middle term whose
filtered multiplicities sum those of the summands. The subobject
and quotient clauses now give crystallinity and additivity of
multiplicities for this second exact sequence. A jump of the
extension contributes a positive multiplicity, and the quotient
has nonnegative multiplicities. Hence some summand has a positive
multiplicity there. The family clause of G2 identifies this
multiplicity with that of the marked fibre, giving its jump
(\nolinkurl{crys_sum}, \nolinkurl{crys_of_inj},
\path{Formalism/AdmissibleDesigns/TensSubsingleton.lean:213--270,273--304};
\nolinkurl{ytG_crys}, \path{Formalism/AdmissibleDesigns/YtRepZero.lean:147--158})
(\path{ytEmbW_of}, \path{TwistedTargetOfInputs.lean:175--198};
\path{ytG_psi}, \path{YtRepZero.lean:107--148};
\path{jump_pi}, \path{ytTarget_of}, \path{YtTargetOf.lean:128--169}).

The row functional on $Y$ extends $O$-linearly to this rational source
with values in $C_e$. Indeed $C_e$ is divisible over a DVR, so a map on
a submodule extends by the injectivity criterion for modules over a
principal ideal domain. On the integral kernel of $p$ its restriction is
the original equivariant map. The orbit difference at $s$ is therefore
sent to the row cocycle. The first clause of the assumed
\path{TBAt} gives $l^{c_B}$ times this cocycle in the finite torsion
condition. For the length, let $U_e$ denote the locally constant torsion
cocycles and $F_e$ the finite torsion condition on representatives.
Transport a representative of the relaxed row quotient to
\[
 \frac{U_e\cap (l^{c_B})^{-1}F_e}{U_e\cap F_e}.
\]
This is well defined and injective: coboundaries are respected by the
coefficient isomorphism, and the second clause of \path{TBAt}
identifies the denominator on cohomology with zero at $e>0$, or with
the unramified classes at $e=0$. In the latter case the group
isomorphism carries inertia onto inertia. Thus membership of the
transported representative in the denominator is exactly membership in
the strict condition. The length of the target is bounded by
$( [E_w:\mathbf Q_l]+1)c_Ba$ by the last clause of \path{TBAt}.
The coefficient ring isomorphism preserves this length, since it
transports submodule chains and $(l)$, and local coefficient
compatibility identifies the degrees
(\path{TBAt}, \path{Formalism/SharpRestriction/LevelLetterCocycles/QuotientHomNonvanishingBound.lean:171--187};
\path{trC_hf}, \path{trC_mem_Z1D}, \path{TrCMemZ1D.lean:146--168};
\path{lenMap_of}, \path{TrivialBoundTargetOfPsi.lean:109--132};
\path{len_gen}, \path{tbLen_w}, \path{TrivialBoundTargetOfDict.lean:77--94,137--159}).

\paragraph{Integral columns and the two lifts.}
Let $T$ be the stable first-block lattice over $O$, $K=\operatorname{Frac}(O)$
and $W=T\otimes_O K$. The map from its integral finite cohomology
$F_f$ to continuous integral $H^1$ in the local model is injective:
an integral primitive there pulls back to an integral primitive here
by the lattice identification and surjectivity of the group isomorphism.
Tate finiteness makes the target finitely generated and $O$ is
Noetherian. To compute the rank, map $F_f$ to the rational finite
cohomology. A continuous rational finite cocycle has compact image,
since the local Galois group is profinite. Dilates of the open lattice
cover its vector space. A finite subcover on the image gives one
nonzero denominator clearing all cocycle values. The multiple still
has a period primitive and hence gives an integral finite class.
Conversely if an integral cocycle becomes $g v-v$ over $K$, clear
one denominator of $v$ to make this an integral coboundary.
Consequently the map is an isomorphism after inverting $O\setminus\{0\}$.
The normalized coordinates identify $W$ with the local first-block
representation $\rho_\pi\chi_1^{-1}$ by a continuous equivariant
linear isomorphism. It conjugates the period coboundary tests in
both directions. This proves $\operatorname{rank}_O(F_f/F_{f,\mathrm{tors}})=j_w$.
Saturation and the basis lifts described above now apply, including
when $j_w=0$. For each of the two evaluations the required lift is
the $l^{m_w}$-multiple lift in the public lattice condition transported
back on cocycles; for the second evaluation its input is the finite
unramified trivial-coefficient class $\delta_O u_w$
(\path{finite_FfG}, \path{bound_compact}, \path{thet_loc},
\path{finrank_Ff}, \path{finFf_of}, \path{MemIntegralSpanIffImage.lean:199--251,265--282,722--850,894--916};
\path{rankFf_normAt}, \path{RankFfNormAt.lean:26--65};
\path{liftCol_eve_Dd}, \path{liftCol_evE_Dd}, \path{FiniteClassesLocalEq.lean:708--786,798--848}).

\paragraph{The off-$Z$ dimension in the determinant.}
At a column and an evaluated point, let $V$ be the full twisted rational
fibre and $W=\rho_\pi\chi_1^{-1}$. The local comparison sends the
finite cocycles $L_x$ and the ordinary coboundaries $\operatorname{im}D$
to their namesakes for $V$, so
$L_x/\operatorname{im}D\simeq H_f^1(V)$ and
$\ker D\simeq H^0(V)$. For the first assertion apply the period
map to the witness and use the inverse coordinate map to recover it.
For the second use $D(v)(g)=gv-v$ and surjectivity onto the local
group. Finite-dimensionality of $H_f^1(V)$ follows by the compact-image
argument for all continuous cocycles and integral Tate finiteness.
Where a stable lattice is needed, take the span of the orbits of an
integral basis lattice; compactness bounds those finitely many basis
orbits in one dilate, so this span is a finite free open lattice
(\path{latSt}, \path{finDim_of_tate}, \path{finX_of},
\path{MemIntegralSpanIffImage.lean:423--469,517--544,968--992};
\path{dimLx_of}, \path{Formalism/AdmissibleDesigns/CohomZeroEquivContinuous.lean:424--488}).

The family condition gives crystallinity of $V$ and equality of its
filtered multiplicities with those at the marked point. The specified
semisimple decomposition there splits off $W$ together with the
characters. The Fontaine direct-sum and subobject clauses give
crystallinity of $W$ and additivity of the multiplicities. All jumps
of these characters are nonnegative by the local gain condition.
Hence $m_V(\tau,i)=m_W(\tau,i)$ for every $i<0$.
For each embedding $\tau$ take a common lower filtration index at which
both exhaustive de Rham filtrations are the whole spaces. Summing
successive dimension differences from that index through $-1$
shows that the dimensions of the two quotients by $\mathrm{Fil}^0$
agree. The chosen coefficient field contains the images of all embeddings of
the local field into the algebraic closure. Restricting the codomain
of each such embedding gives an inverse to composing with the
coefficient-field inclusion. The embedding count therefore gives
the required splitting (\nolinkurl{splits_coef},
\path{Formalism/SharpRestriction/LevelLetterCocycles/LargeLevelLemmas.lean:42--56}).
The Bloch--Kato dimension hypothesis now applies to both crystalline
representations. Their filtered contributions
are equal, and the local gain assumption gives $H^0(W)=0$.
Therefore
\[
 \dim H_f^1(V)=\dim H^0(V)+\dim H_f^1(W)
                =\dim\ker D+j_w .
\]
This proves the required upper bound for the dimension of
$L_x/\operatorname{im}D$ in the determinant argument.
The column comparison consumes this dimension bound through
\nolinkurl{colLx_of_ringDef}, \nolinkurl{colLx} and
\nolinkurl{plugOpen3_B} in \nolinkurl{theorem_n5_of_open_b}
(\path{Formalism/AdmissibleDesigns/ColumnUnramifiedCharacter.lean:69--108};
\path{Formalism/AdmissibleDesigns/ColLx.lean:32--70};
\path{Formalism/AdmissibleDesigns/StatementOneBChangedOfOpen.lean:316--345,351--363}).
The two normalized applications and the public arithmetic hypotheses
are distinguished in Section~\ref{sec:r5-z32-rev}.
These comparisons complete Z3.2 under the hypotheses of
\path{SelmerBoundAtPoint.main}
(\path{crysX_of}, \path{decomp_at}, \path{wp_font}, \path{negX_of},
\path{IsDROfWpIso.lean:83--95,200--270,297--320};
\path{negdim_eq}, \path{count_eq}, \path{rank_le_count}, \path{dimLx_of_font},
\path{Formalism/AdmissibleDesigns/CohomZeroEquivContinuous.lean:199--263,531--600};
\path{open_inputs_at_normalized}, \path{StatementOneBFromUnramified.lean:74--93};
\path{theorem_n5_of_open_b}, \path{StatementOneBChangedOfOpen.lean:351--365}).

Two compatibility statements also enter the cohomology calculation.
The evaluated algebra has its finite-module topology and the twisted
action is continuous for the local coefficient and fraction topologies;
these hypotheses give the locally constant torsion cocycles used in
$\Psi_c$ (\path{topoHyps}, \path{DisplayForms.lean:21--51}).
For the second column lift, $\delta_O u_w$ represents a finite class
by the explicit unramified trivial-coefficient hypothesis of the public
theorem. The conversion is invoked by \nolinkurl{SelmerBoundAtPoint.main}
itself (\nolinkurl{StD.unramified_input_of_admitted},
\path{Formalism/AdmissibleDesigns/UnramifiedInputOfAdmitted.lean:27--34};
\path{Development/AbelianVarietyTheory/Proofs/InertiaEigenvalues/SelmerLengthBoundAtPoint.lean:23--28}). These are additional inputs
to the displayed reduction. The local class field and duality comparison should be read with its
original quantifiers. For a characteristic-zero nonarchimedean local
field, Milne states cup-product duality for finitely generated discrete
modules with torsion of order prime to the characteristic, and the local
Euler characteristic for finite modules of such order. Finite
$l$-primary coefficients over $E_w$ satisfy this characteristic condition
also at $w\mid l$. The passage to the lattice and rational period
conditions uses the local hypotheses and comparisons above
(\path{admBK}, \path{DecompositionData.lean:238--243};
\cite[I.2, Corollary 2.3 and Theorem 2.8]{S-Milne}).
\emph{Status: complete under the hypotheses of \path{SelmerBoundAtPoint.main}, including the explicit local arithmetic, topology, tame and class-field conditions.}

%BIB S-Milne | J. S. Milne | Arithmetic Duality Theorems | second edition, BookSurge, 2006 | https://www.jmilne.org/math/Books/ADTnot.pdf | I.2, Corollary 2.3 and Theorem 2.8, local finite-coefficient duality and Euler characteristic; I.4.10(a),(b) for duality, and the exact sequence following part (c).
\clearpage

